% !TEX program = lualatex
% ═══════════════════════════════════════════════════════════════════════
%  Laboratoire informatique — Introduction à LaTeX
%  Activité sur Overleaf, avec l'aide de l'IA.
%  Jérôme Soucy
%
%  Compilation : lualatex main.tex   (deux passages, pour le sommaire)
%  Les fichiers d'exemple (à distribuer aux étudiants) sont dans
%  exemples/ ; ils se compilent avec pdfLaTeX, comme sur Overleaf.
% ═══════════════════════════════════════════════════════════════════════
\documentclass[aspectratio=169, 11pt]{beamer}
% ═══════════════════════════════════════════════════════════════════════
%  preambule.tex — Réglages, couleurs et raccourcis communs du diaporama
%  du laboratoire « Introduction à LaTeX ».
%  Ce fichier est autonome : il ne dépend d'aucun fichier du livre.
% ═══════════════════════════════════════════════════════════════════════

% ── Encodage, langue, polices (lualatex, comme le livre) ───────────────
\usepackage{fontspec}
\usepackage[french]{babel}

\setmainfont{TeX Gyre Termes}
\IfFontExistsTF{Fira Sans}
  {\setsansfont{Fira Sans}[Scale=MatchLowercase]}
  {\setsansfont{TeX Gyre Heros}[Scale=MatchLowercase]}
\usefonttheme{professionalfonts}
\renewcommand{\familydefault}{\sfdefault}

\usepackage{amsmath}   % avant unicode-math (sinon \underbrace est cassé)
\usepackage{mathtools}
\usepackage{unicode-math}
\IfFontExistsTF{NewCMMath-Book.otf}
  {\setmathfont{NewCMMath-Book}}
  {\setmathfont{Latin Modern Math}}

\usepackage{booktabs}
\usepackage{array}
\usepackage{graphicx}
\usepackage{adjustbox}
\usepackage{siunitx}
\sisetup{output-decimal-marker = {,}, exponent-product = \times,
         group-separator = {\,}, group-minimum-digits = 4}
\usepackage[skins,breakable]{tcolorbox}
\usepackage{tikz}
\usetikzlibrary{calc, arrows.meta, angles, quotes, patterns, positioning,
                decorations.markings, backgrounds, fit, shapes.geometric}

% ── Palette (identique à celle du livre, livrejs.cls) ──────────────────
\definecolor{sectioncolor}{RGB}{28, 63, 110}   % bleu foncé — titres
\definecolor{thmboxbg}{RGB}{228, 238, 252}     % bleu pâle — énoncés
\definecolor{coveraccent}{RGB}{196, 142, 58}   % or chaud — accent
\definecolor{secboxbg}{RGB}{253, 240, 219}     % ocre pâle — « Pour la classe »
\definecolor{steelblue4}{HTML}{36648B}         % figures : objet de départ
\definecolor{firebrick3}{HTML}{CD3333}         % figures : objet image
\definecolor{demobg}{RGB}{247, 247, 247}       % fond neutre — « au tableau »
\definecolor{bgpage}{RGB}{253, 253, 252}       % fond de diapositive, légèrement chaud

% ── Thème beamer : sobre, cohérent avec le livre ────────────────────────
\usecolortheme{default}
\setbeamercolor{palette primary}{bg=sectioncolor, fg=white}
\setbeamercolor{palette secondary}{bg=sectioncolor, fg=white}
\setbeamercolor{palette tertiary}{bg=sectioncolor, fg=white}
\setbeamercolor{title}{fg=sectioncolor}
\setbeamercolor{frametitle}{bg=sectioncolor, fg=white}
\setbeamercolor{structure}{fg=sectioncolor}
\setbeamercolor{section in toc}{fg=sectioncolor}
\setbeamercolor{subsection in toc}{fg=black}
\setbeamercolor{background canvas}{bg=bgpage}
\setbeamercolor{footline}{fg=gray, bg=bgpage}
\setbeamercolor{block title}{bg=sectioncolor, fg=white}
\setbeamercolor{block body}{bg=thmboxbg, fg=black}
\setbeamercolor{item}{fg=sectioncolor}
\setbeamercolor{itemize item}{fg=coveraccent}
\setbeamercolor{itemize subitem}{fg=sectioncolor}
\setbeamertemplate{itemize item}{\raise0.3ex\hbox{\small$\blacktriangleright$}}
\setbeamertemplate{itemize subitem}{\raise0.3ex\hbox{\footnotesize$\bullet$}}
\setbeamertemplate{navigation symbols}{}
\setbeamerfont{frametitle}{series=\bfseries}
\setbeamertemplate{frametitle}{%
  \vskip-1ex
  \begin{beamercolorbox}[wd=\paperwidth,ht=3.8ex,dp=1.3ex,leftskip=1em,rightskip=1em]{frametitle}
    \bfseries\strut\insertframetitle
  \end{beamercolorbox}
  \vskip3pt
}
\setbeamertemplate{footline}{%
  \leavevmode
  \hbox{\begin{beamercolorbox}[wd=\paperwidth,ht=2.4ex,dp=1ex,leftskip=1em,rightskip=1em]{footline}
    {\scriptsize\color{gray}Écrire des mathématiques avec \LaTeX
    \hfill \insertshorttitle{} \hfill \insertframenumber/\inserttotalframenumber}
  \end{beamercolorbox}}%
}
\setbeamertemplate{blocks}[rounded][shadow=false]
\setbeamertemplate{section in toc}{%
  \inserttocsectionnumber.~\inserttocsection\par\vskip0.15em}

% ── Page de section : diapositive de transition ────────────────────────
\AtBeginSection[]{%
  \begin{frame}[plain,c]
    \begin{center}
      {\color{coveraccent}\rule{0.25\textwidth}{1.2pt}}\\[10pt]
      {\usebeamerfont{frametitle}\usebeamercolor[fg]{title}%
       \Large\bfseries\color{sectioncolor}\insertsectionhead}\\[10pt]
      {\color{sectioncolor}\rule{0.45\textwidth}{0.5pt}}
    \end{center}
  \end{frame}
}
\AtBeginSubsection[]{%
  \begin{frame}[plain,c]
    \begin{center}
      {\large\color{gray}\insertsectionhead}\\[4pt]
      {\Large\bfseries\color{coveraccent}\insertsubsectionhead}
    \end{center}
  \end{frame}
}

% ── Raccourcis mathématiques repris du livre (autonomes) ────────────────
\newcommand{\Plan}{\Pi}
\newcommand{\Transf}{\mathcal{T}}
\newcommand{\Aff}{\mathcal{A}}
\newcommand{\Sim}{\mathcal{S}}
\newcommand{\Iso}{\mathcal{I}}
\newcommand{\mat}[4]{\begin{pmatrix}#1 & #2\\ #3 & #4\end{pmatrix}}
\newcommand{\col}[2]{\begin{pmatrix}#1\\ #2\end{pmatrix}}
\newcommand{\coliii}[3]{\begin{pmatrix}#1\\ #2\\ #3\end{pmatrix}}
\newcommand{\norme}[1]{\lVert #1 \rVert}
\newcommand{\Aire}{\operatorname{Aire}}
\newcommand{\dete}{\operatorname{det}}
\renewcommand{\vec}[1]{\mathbf{#1}}
\DeclareMathOperator{\R}{\mathbb{R}}
\DeclareMathOperator{\Rs}{\mathbb{R}^{\star}}
\newcommand{\intff}[2]{\mathopen{[}{#1},{#2}\mathclose{]}}
\newcommand{\intoo}[2]{\mathopen{]}{#1},{#2}\mathclose{[}}
\newcommand{\defi}[1]{{\color{coveraccent}#1}}
\newcommand{\intfo}[2]{\mathopen{[}{#1},{#2}\mathclose{[}}
\DeclareMathOperator{\N}{\mathbb{N}}
\DeclareMathOperator{\Ns}{\mathbb{N}^{\star}}
\DeclareMathOperator{\Q}{\mathbb{Q}}

% ── Raccourcis propres au chapitre 2 (identiques à macros-livre.tex) ────
\newcommand{\dE}{d_E}                         % distance euclidienne
\newcommand{\dT}{d_T}                         % taxi-distance (Manhattan)
\newcommand{\dHam}{d_{\mathrm{Ham}}}          % distance de Hamming
\newcommand{\dH}{d_H}                         % distance de Hausdorff
\newcommand{\boule}[2]{B(#1, #2)}             % boule ouverte
\newcommand{\Vois}[2]{#1^{\oplus #2}}         % epsilon-voisinage (Minkowski)
\newcommand{\Comp}[1]{\mathcal{K}(#1)}        % compacts non vides de E
\newcommand{\comp}[1]{\mathopen{}#1\mathclose{}^{\mathsf{c}}}

% Images bitmap (portraits, figures Python) reprises du livre.
\graphicspath{{exemples/}}

% ── Styles TikZ communs à toutes les figures (identiques au livre) ──────
\tikzset{
  axe/.style        = {-{Stealth[length=2mm]}, gray!70, thin},
  quadrillage/.style= {gray!22, very thin},
  fig/.style        = {fill=steelblue4!22, draw=steelblue4, thick,
                       line join=round},
  figim/.style      = {fill=firebrick3!22, draw=firebrick3, thick,
                       line join=round},
  figaux/.style     = {fill=coveraccent!22, draw=coveraccent!85!black,
                       thick, line join=round},
  pt/.style         = {circle, fill=sectioncolor, inner sep=1.1pt},
  ptim/.style       = {circle, fill=firebrick3, inner sep=1.1pt},
  fleche/.style     = {-{Stealth[length=2.4mm]}, sectioncolor, thick},
  etiq/.style       = {font=\scriptsize, inner sep=1.5pt},
  droite/.style     = {sectioncolor!70, thick, dashed},
}
\newcommand{\repere}[4]{%
  \draw[quadrillage] (#1,#3) grid (#2,#4);
  \draw[axe] (#1,0) -- (#2+0.35,0) node[etiq, below] {$x$};
  \draw[axe] (0,#3) -- (0,#4+0.35) node[etiq, left] {$y$};}
\newcommand{\motifF}{(0,0) -- (0.3,0) -- (0.3,0.55) -- (0.6,0.55) --
  (0.6,0.75) -- (0.3,0.75) -- (0.3,0.85) -- (0.75,0.85) -- (0.75,1) --
  (0,1) -- cycle}

% ── Boîtes d'énoncés (même esprit que le livre : fond bleu pâle) ────────
\tcbset{
  boxenonce/.style={
    enhanced, breakable, colback=thmboxbg, colframe=sectioncolor,
    boxrule=0.6pt, arc=2pt, left=9pt, right=9pt, top=6pt, bottom=6pt,
    fonttitle=\bfseries\sffamily, coltitle=white,
    colbacktitle=sectioncolor, attach boxed title to top left,
    boxed title style={arc=1pt, boxrule=0pt, left=6pt, right=6pt,
      top=3pt, bottom=3pt},
    title style={sharp corners},
  },
  boxclasse/.style={
    enhanced, breakable, colback=secboxbg, colframe=coveraccent,
    boxrule=0.6pt, arc=2pt, left=9pt, right=9pt, top=6pt, bottom=6pt,
    fonttitle=\bfseries\sffamily, coltitle=white,
    colbacktitle=coveraccent, attach boxed title to top left,
    boxed title style={arc=1pt, boxrule=0pt, left=6pt, right=6pt,
      top=3pt, bottom=3pt},
  },
  boxrem/.style={
    enhanced, breakable, colback=white, colframe=steelblue4,
    boxrule=0pt, leftrule=3pt, arc=0pt, left=9pt, right=9pt,
    top=5pt, bottom=5pt,
    colbacktitle=white, coltitle=steelblue4, toptitle=2pt,
    bottomtitle=0pt, fonttitle=\bfseries\sffamily,
  },
  boxoutil/.style={
    enhanced, breakable, colback=white, colframe=steelblue4,
    boxrule=0.6pt, arc=2pt, left=9pt, right=9pt, top=6pt, bottom=6pt,
    fonttitle=\bfseries\sffamily, coltitle=white,
    colbacktitle=steelblue4, attach boxed title to top left,
    boxed title style={arc=1pt, boxrule=0pt, left=6pt, right=6pt,
      top=3pt, bottom=3pt},
  },
  boxtableau/.style={
    enhanced, breakable, colback=demobg, colframe=firebrick3,
    boxrule=0pt, leftrule=4pt, arc=0pt, left=10pt, right=10pt,
    top=7pt, bottom=7pt,
  },
}

% Environnements : titre = « Type N.M — description », fourni en argument.
\newtcolorbox{Enonce}[1]{boxenonce, title={#1}}
\newtcolorbox{Classe}[1]{boxclasse, title={Pour la classe --- #1}}
\newtcolorbox{Outil}[1]{boxoutil, title={Outil --- #1}}
\newtcolorbox{Rem}[1]{boxrem, title={#1}}

% « Démonstration au tableau » : marqueur + rappel très court de l'idée.
\newenvironment{AuTableau}[1][Démonstration au tableau]{%
  \begin{tcolorbox}[boxtableau]
  \bfseries\sffamily\color{firebrick3}\MakeUppercase{#1}\par\nointerlineskip
  \vspace{4pt}\normalfont\rmfamily\upshape\mdseries\color{black}
}{%
  \end{tcolorbox}
}

% Petit renvoi discret vers le numéro du livre, en fin de bloc de titre.
\newcommand{\refnotes}[1]{{\scriptsize\color{gray} (notes de cours, #1)}}

% Insère une figure (tikzpicture ou image) en la contraignant à la fois
% en largeur et en hauteur, sans distorsion : \FigBox[hauteur max]{contenu}.

% ═══════════════════════════════════════════════════════════════════════
%  Environnements propres au laboratoire LaTeX
% ═══════════════════════════════════════════════════════════════════════
\tcbuselibrary{listings}
\definecolor{codebg}{RGB}{246, 248, 251}
\definecolor{vertpiece}{RGB}{58, 125, 68}
\definecolor{promptcol}{RGB}{92, 74, 138}      % violet sobre : les prompts
\definecolor{promptbg}{RGB}{245, 242, 250}

% Code LaTeX : commandes en bleu, commentaires en gris, jamais de gras.
\lstdefinestyle{latexcode}{
  language=[LaTeX]TeX, basicstyle=\ttfamily\scriptsize,
  keywordstyle=\color{sectioncolor}, texcsstyle=*\color{sectioncolor},
  commentstyle=\color{gray}, columns=fullflexible, keepspaces=true,
  upquote=true, showstringspaces=false, breaklines=true,
  moretexcs={documentclass, usepackage, title, author, date, maketitle,
    section, subsection, item, frac, dfrac, sqrt, num, qty, sisetup,
    question, part, parts, choice, CorrectChoice, printanswers, gradetable,
    numpoints, numquestions, pointpoints, hqword, hpword, hsword, htword,
    firstpageheader, runningheader, firstpagefooter, runningfooter,
    fillwithlines, makeemptybox, stretch, vspace, includegraphics, caption,
    centering, text, mathbb, cdot, times, leq, geq, neq, approx, pm, pi,
    infty, sin, cos, log, left, right, vec, quad, textsuperscript, hline,
    draw, node, solutiontitle, renewcommand, newcommand, label, ref, lhead,
    rhead, cfoot, fancyhf, pagestyle, thepage, numpages, textwidth,
    centi, metre, squared, minute, alpha, beta, theta, Delta, in, to,
    displaystyle, ldots, thispagestyle, tikz},
  literate={é}{{é}}1 {è}{{è}}1 {ê}{{ê}}1 {à}{{à}}1 {â}{{â}}1
           {ç}{{ç}}1 {ù}{{ù}}1 {û}{{û}}1 {ô}{{ô}}1 {î}{{î}}1 {ï}{{ï}}1
           {É}{{É}}1 {À}{{À}}1 {«}{{«}}1 {»}{{»}}1 {œ}{{œ}}1 {’}{{’}}1
           {—}{{—}}1 {–}{{–}}1
}

\tcbset{
  codestyle/.style={
    enhanced, colback=codebg, colframe=sectioncolor!35, boxrule=0.4pt,
    arc=2pt, left=6pt, right=6pt, top=1pt, bottom=1pt,
    listing options={style=latexcode},
    fonttitle=\sffamily\footnotesize, coltitle=white,
    colbacktitle=sectioncolor!80, attach boxed title to top left={yshift=-2mm, xshift=3mm},
    boxed title style={arc=1pt, boxrule=0pt, left=4pt, right=4pt, top=1pt, bottom=1pt},
  },
}
% Code seul : \begin{CodeTeX}[title=...] ... \end{CodeTeX}
\newtcblisting{CodeTeX}[1][]{codestyle, listing only, #1}
% Code et rendu côte à côte (le code est réellement compilé dans la diapo).
\newtcblisting{CodeRendu}[1][]{codestyle, listing side text,
  lefthand ratio=0.56, sidebyside gap=4mm, sidebyside align=center seam,
  segmentation style={draw=sectioncolor!35, dashed},
  fontlower=\footnotesize, #1}

% Prompt type pour l'IA : bulle de conversation sobre.
\newcommand{\iconebulle}{%
  \tikz[baseline=-0.7ex]{\fill[white, rounded corners=1pt] (0,-0.08) rectangle (0.34,0.16);
    \fill[white] (0.06,-0.07) -- (0.02,-0.17) -- (0.14,-0.07) -- cycle;}}
\newtcolorbox{Prompt}[1]{
  enhanced, colback=promptbg, colframe=promptcol, boxrule=0pt,
  leftrule=3pt, arc=0pt, outer arc=0pt, left=9pt, right=9pt, top=5pt,
  bottom=5pt, fontupper=\footnotesize\sffamily\raggedright,
  fonttitle=\sffamily\footnotesize, coltitle=white, colbacktitle=promptcol,
  title={\iconebulle\hspace{0.45em}Prompt type --- #1},
  attach boxed title to top left={yshift=-2mm, xshift=4mm},
  boxed title style={arc=1pt, boxrule=0pt, left=5pt, right=5pt, top=2pt, bottom=2pt},
  before upper={\setlength{\parskip}{3pt}},
}
% Un champ à compléter dans un prompt : [niveau], [sujet], etc.
\newcommand{\champ}[1]{{\color{promptcol}[#1]}}
% Une commande LaTeX citée dans le texte : \cmd{section} → \section
\newcommand{\cmd}[1]{\texttt{\textbackslash #1}}

% Exercice en laboratoire : même esprit que « Pour la classe ».
\newtcolorbox{Exercice}[1]{boxclasse, title={À vous --- #1}}

\newcommand{\FigBox}[2][0.56\textheight]{%
  \begin{center}
    \adjustbox{max width=\linewidth, max height=#1}{#2}
  \end{center}}


\title[Laboratoire LaTeX]{Écrire des mathématiques avec \LaTeX}
\subtitle{Laboratoire informatique --- Formation à l'enseignement des mathématiques}
\author{Jérôme Soucy}
\institute{Université Laval --- Département de mathématiques et de statistique}
\date{}

\begin{document}

\begin{frame}[plain]
  \begin{center}
    {\color{sectioncolor}\rule{0.5\textwidth}{0.8pt}}\\[3pt]
    {\color{coveraccent}\rule{0.14\textwidth}{1.8pt}}

    \vspace{0.55cm}
    {\large\color{gray}Laboratoire informatique}\\[3pt]
    {\LARGE\bfseries\color{sectioncolor}Écrire des mathématiques avec \LaTeX}

    \vspace{0.35cm}
    {\normalsize\color{sectioncolor}Overleaf, l'IA, et ce qu'il faut comprendre}

    \vspace{0.5cm}
    {\color{sectioncolor}%
      $\displaystyle x = \frac{-b \pm \sqrt{b^2 - 4ac}}{2a}$}

    \vspace{0.5cm}
    {\sffamily\color{sectioncolor}Jérôme Soucy}\\[1pt]
    {\small\color{gray}Université Laval}

    \vspace{0.35cm}
    {\color{coveraccent}\rule{0.14\textwidth}{1.8pt}}\\[3pt]
    {\color{sectioncolor}\rule{0.5\textwidth}{0.8pt}}
  \end{center}
\end{frame}

\begin{frame}{Le déroulement du laboratoire}
  \small
  \begin{center}
  \begin{tabular}{rll}
    \toprule
    & Étape & Ce que vous faites \\
    \midrule
    1 & Pourquoi \LaTeX ? & on écoute, on discute \\
    2 & Overleaf & créer un compte, un premier projet \\
    3 & Anatomie d'un document & modifier un fichier existant \\
    4 & Compiler et corriger & réparer des erreurs \\
    5 & Le mode mathématique & reproduire des formules \\
    6 & Travailler avec l'IA & écrire de bons prompts \\
    7 & Livrable 1 : une activité & fiche de 2 à 3 pages \\
    8 & Livrable 2 : un examen & classe \texttt{exam}, corrigé \\
    9 & Remise & vérifier, déposer les PDF \\
    \bottomrule
  \end{tabular}
  \end{center}
  \vspace{4pt}
  Trois fichiers d'exemple vous sont fournis : un document minimal, une
  fiche d'activité et un examen. Ils compilent tels quels sur Overleaf.
\end{frame}

\begin{frame}{Plan de l'exposé}
  \footnotesize
  \begin{columns}[T]
    \begin{column}{0.5\textwidth}
      \tableofcontents[hideallsubsections,sections={1-5}]
    \end{column}
    \begin{column}{0.5\textwidth}
      \tableofcontents[hideallsubsections,sections={6-10}]
    \end{column}
  \end{columns}
\end{frame}

% ═══════════════════════════════════════════════════════════════════════
\section{Pourquoi \texorpdfstring{\LaTeX}{LaTeX} ?}
% ═══════════════════════════════════════════════════════════════════════

\begin{frame}{Le standard de l'écriture mathématique}
  \small
  \begin{columns}[c]
    \begin{column}{0.5\textwidth}
      \begin{itemize}
        \item Créé par Donald Knuth (\TeX, 1978) puis Leslie Lamport
          (\LaTeX, 1984), pour composer des mathématiques avec la qualité
          d'un livre imprimé.
        \item Utilisé pour presque tous les articles, manuels et notes de
          cours de mathématiques --- y compris les notes de ce cours.
        \item Gratuit, ouvert, et le même fichier compile partout, dans
          dix ans comme aujourd'hui.
      \end{itemize}
    \end{column}
    \begin{column}{0.46\textwidth}
      \begin{tcolorbox}[codestyle]
        \centering
        $\displaystyle \sum_{k=1}^{n} k^2 = \frac{n(n+1)(2n+1)}{6}$
        \\[8pt]
        $\displaystyle \int_0^{\pi} \sin x \,\mathrm{d}x = 2$
        \\[8pt]
        $\displaystyle \lim_{n\to\infty}\left(1+\frac1n\right)^{\!n} = \mathrm{e}$
      \end{tcolorbox}
      \centering\scriptsize Trois lignes de code chacune, sans souris.
    \end{column}
  \end{columns}
\end{frame}

\begin{frame}{Traitement de texte ou \LaTeX ?}
  \small
  \begin{center}
  \begin{tabular}{p{3.2cm}p{4.3cm}p{4.3cm}}
    \toprule
    & Traitement de texte & \LaTeX \\
    \midrule
    Formules & éditeur d'équations, à la souris & tapées au clavier, rendu
      de qualité livre \\
    Mise en page & vous la faites, à la main & calculée pour vous,
      cohérente partout \\
    Numérotation & à corriger quand on insère & automatique (sections,
      questions, figures) \\
    Format du fichier & binaire & texte brut : lisible, réutilisable,
      généré par l'IA \\
    Ce qu'on voit & le résultat, en direct & le code, puis le PDF
      compilé \\
    \bottomrule
  \end{tabular}
  \end{center}
  \vspace{4pt}
  On ne dessine pas la page : on décrit la structure (un titre, une
  question, une formule) et \LaTeX\ s'occupe de la mise en page.
\end{frame}

\begin{frame}{Pourquoi c'est particulièrement utile en enseignement}
  \begin{Classe}{Ce que \LaTeX\ change pour un enseignant de mathématiques}
    \small
    \begin{itemize}
      \item Un examen et son corrigé dans un seul fichier : un
        interrupteur affiche ou cache les solutions.
      \item Le barème se calcule tout seul : points par question,
        total, grille de correction.
      \item Des banques de questions réutilisables d'une année à l'autre,
        d'un groupe à l'autre (version A, version B).
      \item Des figures précises (TikZ) : triangles, repères, graphes de
        fonctions, à la bonne échelle.
      \item Des documents cohérents pour toute une équipe (école ou département) :
        mêmes en-têtes, même notation, même mise en page.
      \item Avec l'IA, le premier jet ne prend que quelques minutes ;
        le travail d'enseignant, c'est d'en vérifier et d'en ajuster le
        contenu.
    \end{itemize}
  \end{Classe}
\end{frame}

\begin{frame}{Ce qu'il faut accepter}
  \small
  \begin{itemize}
    \item \LaTeX\ n'est pas un logiciel « ce que vous voyez est ce que vous
      obtenez » : on écrit du code, on compile, on regarde le PDF.
    \item Une petite faute de frappe (une accolade oubliée) peut empêcher
      la compilation.
    \item L'IA écrit très bien du \LaTeX, mais elle se trompe aussi :
      commandes inventées, extensions inutiles, mathématiques fausses,
      texte en anglais.
  \end{itemize}
  \vspace{4pt}
  \begin{Rem}{L'objectif d'aujourd'hui}
    Pas de mémoriser des centaines de commandes. Comprendre assez la
    structure d'un fichier pour lire ce que l'IA produit, le modifier à la
    main sans tout casser, et réparer une erreur de compilation.
  \end{Rem}
\end{frame}

% ═══════════════════════════════════════════════════════════════════════
\section{Overleaf : \texorpdfstring{\LaTeX}{LaTeX} dans le navigateur}
% ═══════════════════════════════════════════════════════════════════════

\begin{frame}{Démarrer sur Overleaf}
  \small
  \begin{enumerate}
    \item Aller sur \texttt{overleaf.com} et créer un compte gratuit.
    \item Nouveau projet (New Project) $\to$ Projet vide (Blank Project),
      lui donner un nom.
    \item Pour partir d'un fichier fourni : Nouveau projet $\to$
      Téléverser un projet (Upload Project), ou téléverser le
      \texttt{.tex} dans un projet existant.
    \item Rien à installer : la compilation se fait sur les serveurs
      d'Overleaf.
  \end{enumerate}
  \vspace{4pt}
  \begin{Rem}{Bon à savoir}
    \small Le compilateur par défaut est pdfLaTeX : c'est celui que
    supposent nos fichiers d'exemple. On peut le changer dans le
    menu (Menu $\to$ Compilateur), mais ce n'est pas nécessaire
    aujourd'hui.
  \end{Rem}
\end{frame}

\begin{frame}{L'interface}
  \FigBox[0.6\textheight]{%
    \begin{tikzpicture}[x=1cm, y=1cm, font=\scriptsize\sffamily]
      \draw[sectioncolor, thick, rounded corners=2pt] (0,0) rectangle (13,6.2);
      \fill[sectioncolor!12] (0,5.6) rectangle (13,6.2);
      \node[anchor=west] at (0.2,5.9) {Menu \quad Mon projet};
      \node[fill=vertpiece!80!black, text=white, rounded corners=2pt,
        inner sep=3pt] at (8.3,5.9) {Recompiler};
      \draw[sectioncolor!60] (2.6,0) -- (2.6,5.6);
      \draw[sectioncolor!60] (7.4,0) -- (7.4,5.6);
      \node[anchor=north west, align=left] at (0.15,5.45)
        {Fichiers\\[2pt] \texttt{main.tex}\\ \texttt{figure.png}};
      \node[anchor=north west, align=left, font=\scriptsize\ttfamily] at (2.75,5.45)
        {\textbackslash documentclass\{article\}\\
         \textbackslash usepackage[french]\{babel\}\\[2pt]
         \textbackslash begin\{document\}\\
         Bonjour !\\
         \textbackslash end\{document\}};
      \fill[white] (7.7,0.3) rectangle (12.7,5.3);
      \draw[gray!60] (7.7,0.3) rectangle (12.7,5.3);
      \node at (10.2,4.6) {Bonjour !};
      \node[firebrick3, align=center] at (1.3,-0.5) {1. les fichiers\\ du projet};
      \node[firebrick3, align=center] at (5,-0.5) {2. l'éditeur :\\ le code \texttt{.tex}};
      \node[firebrick3, align=center] at (10.2,-0.5) {3. l'aperçu :\\ le PDF compilé};
      \node[firebrick3, align=center] at (10.4,6.65) {4. compiler (Ctrl+Entrée,\\ Cmd+Entrée sur Mac)};
    \end{tikzpicture}}
  \small On modifie le code à gauche, on recompile, on regarde le PDF à
  droite. Les erreurs s'affichent près du bouton Recompiler.
\end{frame}

\begin{frame}{Le cycle de travail}
  \FigBox[0.32\textheight]{%
    \begin{tikzpicture}[
        etape/.style={draw=sectioncolor, fill=thmboxbg, rounded corners=2pt,
          thick, align=center, font=\small, minimum height=1.1cm,
          text width=2.8cm}]
      \node[etape] (a) at (0,0) {Écrire ou modifier\\le code \texttt{.tex}};
      \node[etape] (b) at (4.2,0) {Compiler\\(pdfLaTeX)};
      \node[etape] (c) at (8.4,0) {Lire le PDF\\(ou l'erreur)};
      \draw[fleche] (a) -- (b);
      \draw[fleche] (b) -- (c);
      \draw[fleche] (c.south) to[out=-150, in=-30] (a.south);
    \end{tikzpicture}}
  \small
  \begin{itemize}
    \item Compiler souvent : après chaque petite modification.
    \item Une erreur à la fois : la première erreur signalée est presque
      toujours la vraie ; les suivantes en découlent souvent.
    \item Le PDF reste celui de la dernière compilation réussie : vérifier
      qu'il a bien changé.
  \end{itemize}
\end{frame}

\begin{frame}{À vous}
  \begin{Exercice}{Un premier projet}
    \small
    \begin{enumerate}
      \item Créez un projet sur Overleaf et téléversez le fichier
        \texttt{01-minimal.tex}.
      \item Compilez-le : vous devez obtenir une page avec trois sections.
      \item Remplacez « Prénom Nom » par votre nom, recompilez.
      \item Ajoutez un quatrième élément à la liste, recompilez.
    \end{enumerate}
  \end{Exercice}
\end{frame}

% ═══════════════════════════════════════════════════════════════════════
\section{L'anatomie d'un document}
% ═══════════════════════════════════════════════════════════════════════

\begin{frame}[fragile]{Trois parties, toujours dans cet ordre}
  \begin{columns}[c]
    \begin{column}{0.6\textwidth}
\begin{CodeTeX}[title={main.tex}]
\documentclass[12pt]{article}   % 1. la classe

% 2. le préambule : les réglages
\usepackage[T1]{fontenc}
\usepackage[french]{babel}
\usepackage{amsmath, amssymb}
\title{Mon premier document}
\author{Prénom Nom}

\begin{document}               % 3. le corps
\maketitle
Le texte et les formules vont ici.
\end{document}
\end{CodeTeX}
    \end{column}
    \begin{column}{0.37\textwidth}
      \small
      \begin{enumerate}
        \item La classe : le type de document.
        \item Le préambule : extensions et réglages. Rien ne s'y
          affiche.
        \item Le corps, entre \cmd{begin\{document\}} et
          \cmd{end\{document\}} : tout ce qui sera imprimé.
      \end{enumerate}
      \vspace{4pt}
      Tout ce qui suit un \texttt{\%} sur une ligne est un commentaire,
      ignoré à la compilation.
    \end{column}
  \end{columns}
\end{frame}

\begin{frame}{Les classes de documents}
  \small \texttt{\textbackslash documentclass[options]\{classe\}} : la
  première ligne décide du type de document.
  \begin{center}
  \begin{tabular}{lll}
    \toprule
    Classe & Pour quoi faire & Exemple en enseignement \\
    \midrule
    \texttt{article} & document court, sections & fiche d'activité,
      plan de leçon \\
    \texttt{report}, \texttt{book} & document long, chapitres & notes de
      cours, cahier d'exercices \\
    \texttt{exam} & questions, points, solutions & examen, quiz,
      devoir \\
    \texttt{beamer} & diapositives & présentation en classe (comme
      celle-ci) \\
    \bottomrule
  \end{tabular}
  \end{center}
  \vspace{4pt}
  Options courantes : taille du texte (\texttt{10pt}, \texttt{11pt},
  \texttt{12pt}), format du papier (\texttt{letterpaper} est le défaut en
  Amérique du Nord), \texttt{twocolumn}.
\end{frame}

\begin{frame}[fragile]{Un préambule en français, et pas plus}
  \begin{columns}[c]
    \begin{column}{0.56\textwidth}
\begin{CodeTeX}[title={Le préambule de nos exemples}]
\usepackage[T1]{fontenc}
\usepackage[french]{babel}
\usepackage{amsmath, amssymb}
\usepackage{siunitx}
\sisetup{output-decimal-marker = {,}}
\usepackage[margin=2.5cm]{geometry}
\usepackage{graphicx}
\usepackage{tikz}
\end{CodeTeX}
    \end{column}
    \begin{column}{0.41\textwidth}
      \footnotesize
      \begin{itemize}
        \item \texttt{fontenc}, \texttt{babel} : accents, césure et
          typographie françaises (espaces avant « : ; ? ! »).
        \item \texttt{amsmath}, \texttt{amssymb} : outils et symboles
          mathématiques.
        \item \texttt{siunitx} : nombres et unités, virgule décimale.
        \item \texttt{geometry} : les marges.
        \item \texttt{graphicx} : insérer des images.
        \item \texttt{tikz} : dessiner des figures.
      \end{itemize}
    \end{column}
  \end{columns}
  \vspace{4pt}
  \begin{Rem}{Conseil}
    \small Un préambule court est un préambule qui compile. Si l'IA vous
    propose quinze extensions, demandez-lui à quoi sert chacune.
  \end{Rem}
\end{frame}

\begin{frame}[fragile]{Commandes et environnements}
  \small Deux formes à reconnaître, partout :
\begin{CodeTeX}[title={Une commande, avec option et argument}]
\commande[option facultative]{argument obligatoire}
\section{Mise en situation}       \usepackage[french]{babel}
\end{CodeTeX}
\begin{CodeTeX}[title={Un environnement : un début et une fin qui se répondent}]
\begin{nom}
   ...contenu...
\end{nom}
\end{CodeTeX}
  \begin{itemize}
    \item Les accolades \texttt{\{ \}} entourent toujours un argument ; elles
      vont par paires.
    \item Les majuscules comptent : \cmd{Section} n'existe pas.
    \item Chaque \cmd{begin\{x\}} doit avoir son \cmd{end\{x\}}.
  \end{itemize}
\end{frame}

\begin{frame}[fragile]{Structurer le texte}
  \begin{columns}[T]
    \begin{column}{0.55\textwidth}
\begin{CodeTeX}[title={Sections et paragraphes}]
\section{Mise en situation}
Un paragraphe se termine par
une ligne vide.

Un nouveau paragraphe.

\subsection{Matériel}
Calculatrice, règle.
\end{CodeTeX}
    \end{column}
    \begin{column}{0.42\textwidth}
      \small
      \begin{itemize}
        \item La numérotation des sections est automatique.
        \item \cmd{section*\{...\}} : une section non numérotée.
        \item Un seul retour à la ligne ne fait rien ; une ligne vide
          commence un paragraphe.
        \item Les espaces multiples comptent pour un seul.
      \end{itemize}
    \end{column}
  \end{columns}
\end{frame}

\begin{frame}[fragile]{Les listes}
\begin{CodeRendu}[title={Numérotée ou à puces}]
\begin{enumerate}
  \item Lire l'énoncé.
  \item Faire un schéma.
\end{enumerate}
\begin{itemize}
  \item une puce
\end{itemize}
\end{CodeRendu}
  \small Chaque élément commence par \cmd{item}. Dans un document de classe
  \texttt{article}, le rendu est plus sobre que dans ces diapositives.
\end{frame}

\begin{frame}{Les caractères réservés}
  \small Dix caractères ont un sens particulier pour \LaTeX. Pour les
  afficher :
  \begin{center}
  \begin{tabular}{lll}
    \toprule
    Caractère & Son rôle & Pour l'écrire \\
    \midrule
    \texttt{\%} & commentaire & \texttt{\textbackslash\%} \\
    \texttt{\$} & mode mathématique & \texttt{\textbackslash\$} \\
    \texttt{\&} & colonnes d'un tableau & \texttt{\textbackslash\&} \\
    \texttt{\#} & paramètre d'une commande & \texttt{\textbackslash\#} \\
    \texttt{\_} \quad \texttt{\^{}} & indice, exposant (en mode math) &
      \texttt{\textbackslash\_} \quad \texttt{\textbackslash\^{}\{\}} \\
    \texttt{\{} \quad \texttt{\}} & arguments & \texttt{\textbackslash\{}
      \quad \texttt{\textbackslash\}} \\
    \texttt{\textbackslash} & début d'une commande &
      \texttt{\textbackslash textbackslash} \\
    \texttt{\~{}} & espace insécable & \texttt{\textbackslash\~{}\{\}} \\
    \bottomrule
  \end{tabular}
  \end{center}
  \begin{AuTableau}[Piège fréquent]
    \small « Une hausse de 15\% » : sans la barre oblique, \LaTeX\
    ignore la fin de la ligne. Le symbole \texttt{\%} est le plus souvent
    oublié.
  \end{AuTableau}
\end{frame}

\begin{frame}{À vous}
  \begin{Exercice}{Modifier à la main}
    \small Dans \texttt{01-minimal.tex} :
    \begin{enumerate}
      \item Passez de \texttt{12pt} à \texttt{11pt} ; observez le PDF.
      \item Ajoutez une section « Conclusion » avec deux paragraphes.
      \item Écrivez la phrase « Le rabais est de 20 \% sur 45 \$. » de
        façon qu'elle s'affiche correctement.
      \item Transformez la liste numérotée en liste à puces.
    \end{enumerate}
  \end{Exercice}
\end{frame}

% ═══════════════════════════════════════════════════════════════════════
\section{Compiler et corriger les erreurs}
% ═══════════════════════════════════════════════════════════════════════

\begin{frame}{Ce qui se passe quand on compile}
  \small
  \begin{itemize}
    \item Le compilateur (pdfLaTeX sur Overleaf) lit le fichier
      \texttt{.tex} de haut en bas, exécute les commandes et produit un
      PDF.
    \item Il écrit aussi un journal (le « log ») : avertissements et
      erreurs, avec le numéro de la ligne en cause.
    \item Avertissement (jaune) : le PDF est produit, mais quelque chose
      est à vérifier (une ligne trop longue, une référence inconnue).
    \item Erreur (rouge) : la compilation s'est arrêtée ou a dû deviner.
      Il faut la corriger, même si un PDF s'affiche.
  \end{itemize}
  \vspace{4pt}
  \begin{Rem}{Sur Overleaf}
    \small Le bouton à côté de Recompiler (Journaux et fichiers produits,
    ou Logs and output files) liste les problèmes ; un clic sur une erreur
    amène à la ligne correspondante dans l'éditeur.
  \end{Rem}
\end{frame}

\begin{frame}{Les quatre erreurs que vous verrez le plus}
  \footnotesize
  \begin{center}
  \begin{tabular}{p{3.6cm}p{4.2cm}p{4.9cm}}
    \toprule
    Message & Cause probable & Correction \\
    \midrule
    \texttt{Undefined control sequence} & une commande mal écrite
      (\cmd{fracc}) ou dont l'extension n'est pas chargée & corriger
      l'orthographe, ou ajouter le \cmd{usepackage} \\
    \texttt{Missing \$ inserted} & des mathématiques hors du mode
      mathématique (\texttt{x\^{}2} dans le texte, un \texttt{\_}) &
      entourer de \texttt{\$...\$} \\
    \texttt{Runaway argument}, \texttt{Missing \} inserted} & une
      accolade oubliée & compter les \texttt{\{} et les \texttt{\}} \\
    \texttt{\textbackslash begin\{x\} ended by \textbackslash end\{y\}} &
      environnements mal fermés ou croisés & refermer dans l'ordre
      inverse d'ouverture \\
    \texttt{File `xyz.sty' not found} & extension inexistante (souvent
      inventée par l'IA) & retirer la ligne ou demander une
      alternative \\
    \bottomrule
  \end{tabular}
  \end{center}
  \small Le numéro de ligne s'écrit \texttt{l.12} : c'est là que \LaTeX\
  s'est aperçu du problème --- la cause est à cet endroit ou juste avant.
\end{frame}

\begin{frame}{Une méthode pour ne pas paniquer}
  \small
  \begin{enumerate}
    \item Lire la première erreur seulement, et sa ligne.
    \item Regarder cette ligne et celle d'avant : accolade, \texttt{\$},
      orthographe, \cmd{end} manquant ?
    \item Si on ne trouve pas : mettre des \texttt{\%} devant le bloc
      suspect pour le neutraliser, recompiler, réduire le bloc.
    \item Toujours pas ? Demander à l'IA en lui donnant le message
      d'erreur complet et les lignes concernées (prompt plus loin).
  \end{enumerate}
  \vspace{4pt}
  \begin{AuTableau}[Le bon réflexe]
    \small Compiler après chaque petite modification. Une erreur dans
    trois lignes que l'on vient d'écrire se trouve en dix secondes ; dans
    trois pages, elle peut prendre une demi-heure.
  \end{AuTableau}
\end{frame}

\begin{frame}[fragile]{À vous}
  \begin{columns}[c]
    \begin{column}{0.52\textwidth}
      \begin{Exercice}{Réparer}
        \small Copiez ces lignes dans le corps de votre document,
        compilez, lisez les messages, et corrigez les cinq erreurs.
      \end{Exercice}
    \end{column}
    \begin{column}{0.45\textwidth}
\begin{CodeTeX}[title={Cinq erreurs cachées}]
\section{Exercice
Calculer x^2 si $x = 3$.
\begin{itemize}
  \item La fraction $\frac{1}{2$.
  \item Le rabais : 20%.
\end{enumerate}
\end{CodeTeX}
    \end{column}
  \end{columns}
\end{frame}

% ═══════════════════════════════════════════════════════════════════════
\section{Le mode mathématique}
% ═══════════════════════════════════════════════════════════════════════

\begin{frame}[fragile]{Dans le texte, ou hors du texte}
\begin{CodeRendu}[title={Trois façons d'entrer en mode mathématique}]
Le côté $c$ vérifie $a^2 + b^2 = c^2$.

Hors texte, centré :
\[ c = \sqrt{a^2 + b^2} \]

Numéroté :
\begin{equation}
  A = \pi r^2
\end{equation}
\end{CodeRendu}
  \small
  \begin{itemize}
    \item \texttt{\$...\$} : une formule dans la phrase.
    \item \texttt{\textbackslash[ ... \textbackslash]} : une formule importante, sur sa
      propre ligne.
    \item \texttt{equation} : la même chose, avec un numéro qu'on peut citer.
  \end{itemize}
\end{frame}

\begin{frame}[fragile]{L'essentiel (1) : exposants, fractions, racines}
\begin{CodeRendu}[title={Exposants et indices}]
$x^2 \quad x^{10} \quad a_1
 \quad a_{n+1} \quad 2^{-3}$
\end{CodeRendu}
\begin{CodeRendu}[title={Fractions et racines}]
$\frac{3}{4} \quad \dfrac{x+1}{x-2}
 \quad \sqrt{2} \quad \sqrt[3]{27}$
\end{CodeRendu}
  \small Plus d'un caractère en exposant ou en indice : des accolades
  (\texttt{x\^{}10} donne $x^10$, et non $x^{10}$). \cmd{dfrac} : une
  fraction en grand, même dans le texte.
\end{frame}

\begin{frame}[fragile]{L'essentiel (2) : opérations et symboles}
\begin{CodeRendu}[title={Opérations et relations}]
$3 \times 4 \quad 2 \cdot x \quad a \div b$

$x \leq 5 \quad y \geq 0 \quad a \neq b
 \quad \pi \approx 3{,}14$
\end{CodeRendu}
\begin{CodeRendu}[title={Lettres grecques et ensembles}]
$\alpha \quad \beta \quad \theta
 \quad \Delta \quad \pi$

$x \in \mathbb{R} \quad n \in \mathbb{N}$
\end{CodeRendu}
\end{frame}

\begin{frame}[fragile]{L'essentiel (3) : fonctions et texte}
\begin{CodeRendu}[title={Fonctions : droites et bien espacées}]
$\sin(30^\circ) = \frac12
 \quad \log_2 8 = 3 \quad \cos \theta$
\end{CodeRendu}
\begin{CodeRendu}[title={Du texte dans une formule}]
$P = 2(L + l) \text{ où } L
 \text{ est la longueur}$
\end{CodeRendu}
  \small En mode mathématique, les espaces sont ignorés et les lettres
  sont en italique : on écrit \cmd{sin}, jamais \texttt{sin}, et
  \cmd{text\{...\}} pour des mots.
\end{frame}

\begin{frame}[fragile]{Parenthèses et systèmes}
\begin{CodeRendu}[title={Parenthèses qui s'ajustent}]
\[ \left( \frac{a}{b} \right)^2
   = \frac{a^2}{b^2} \]
\end{CodeRendu}
\begin{CodeRendu}[title={Un système d'équations}]
\[ \begin{cases}
     2x + y = 7 \\
     x - y = 2
   \end{cases} \]
\end{CodeRendu}
\end{frame}

\begin{frame}[fragile]{Une démarche alignée}
\begin{CodeRendu}[title={L'environnement align*}]
\begin{align*}
  3x - 5 &= 10 \\
  3x &= 15 \\
  x &= 5
\end{align*}
\end{CodeRendu}
  \small
  \begin{itemize}
    \item \texttt{\&} marque le point d'alignement (ici, avant le signe
      $=$) ; \texttt{\textbackslash\textbackslash} passe à la ligne suivante.
    \item Sans l'étoile, \texttt{align} numérote chaque ligne.
    \item Idéal pour un corrigé : une étape par ligne, toutes alignées.
  \end{itemize}
\end{frame}

\begin{frame}[fragile]{Les nombres à la française}
  \begin{AuTableau}[Piège fréquent]
    \footnotesize En mode mathématique, la virgule est une ponctuation :
    \LaTeX\ ajoute une espace après elle, et $3,14$ s'affiche avec un trou.
    Avec \texttt{siunitx} réglé sur la virgule décimale, \cmd{num} écrit
    les décimaux et groupe les chiffres ; \cmd{qty} ajoute l'unité.
  \end{AuTableau}
\begin{CodeRendu}[title={Trois solutions}]
Mauvais : $3,14$

Accolades : $3{,}14$

Avec siunitx : $\num{3,14}$
et $\num{1234567,89}$

Unités : \qty{12,5}{\centi\metre},
\qty{4}{\centi\metre\squared}
\end{CodeRendu}
\end{frame}

\begin{frame}[fragile]{Trouver un symbole qu'on ne connaît pas}
  \small
  \begin{itemize}
    \item Detexify (\texttt{detexify.kirelabs.org}) : on dessine le
      symbole, le site propose la commande.
    \item L'IA : « Quelle est la commande \LaTeX\ du symbole “environ
      égal” ? »
    \item La documentation d'Overleaf (\texttt{overleaf.com/learn}).
  \end{itemize}
\begin{CodeRendu}[title={Quelques symboles utiles en géométrie}]
$\angle ABC \quad 45^\circ \quad \overline{AB}
 \quad \vec{u} \quad \parallel \quad \perp
 \quad \sim \quad \cong$
\end{CodeRendu}
\end{frame}

\begin{frame}{À vous}
  \begin{Exercice}{Reproduire ces formules}
    \small Dans votre document, écrivez le code qui produit exactement :
    \begin{enumerate}
      \item $f(x) = -2(x - 3)^2 + 1$
      \item L'aire d'un disque de rayon \qty{5}{\centi\metre} est
        $A = \pi r^2 \approx \qty{78,54}{\centi\metre\squared}$.
      \item \( \displaystyle x = \frac{-b \pm \sqrt{b^2 - 4ac}}{2a} \),
        hors texte et centrée.
      \item La démarche $2(x + 4) = 18$, puis $x + 4 = 9$, puis $x = 5$,
        écrite sur trois lignes alignées sur les signes $=$.
      \item $\sin(\theta) = \dfrac{\text{opposé}}{\text{hypoténuse}}$
    \end{enumerate}
    Les solutions sont à la fin de la présentation.
  \end{Exercice}
\end{frame}

% ═══════════════════════════════════════════════════════════════════════
\section{Images, tableaux et figures}
% ═══════════════════════════════════════════════════════════════════════

\begin{frame}[fragile]{Insérer une image}
  \begin{columns}[c]
    \begin{column}{0.58\textwidth}
\begin{CodeTeX}[title={Avec graphicx}]
\begin{figure}[h]
  \centering
  \includegraphics[width=0.5\textwidth]{triangle.png}
  \caption{Un triangle rectangle.}
\end{figure}
\end{CodeTeX}
    \end{column}
    \begin{column}{0.39\textwidth}
      \small
      \begin{enumerate}
        \item Téléverser l'image dans le projet Overleaf (icône
          Téléverser, au-dessus de la liste des fichiers).
        \item Écrire son nom exact, avec l'extension.
        \item Régler la taille par \texttt{width=}, en fraction de la
          largeur du texte.
      \end{enumerate}
    \end{column}
  \end{columns}
  \vspace{4pt}
  \small \texttt{figure} numérote l'image et la place au mieux dans la
  page ; \texttt{[h]} demande « ici, si possible ».
\end{frame}

\begin{frame}[fragile]{Un tableau de valeurs}
\begin{CodeRendu}[title={L'environnement tabular}]
\begin{tabular}{c|ccccc}
  $x$    & $-2$ & $-1$ & $0$ & $1$ & $2$ \\
  \hline
  $f(x)$ & $4$  & $1$  & $0$ & $1$ & $4$
\end{tabular}
\end{CodeRendu}
  \small
  \begin{itemize}
    \item \texttt{\{c|ccccc\}} : une colonne par lettre (\texttt{l}, \texttt{c},
      \texttt{r} : à gauche, centrée, à droite) ; \texttt{|} trace un
      trait vertical.
    \item \texttt{\&} sépare les colonnes, \texttt{\textbackslash\textbackslash} termine
      une ligne, \cmd{hline} trace un trait horizontal.
    \item Pour un grand tableau, l'IA fait gagner beaucoup de temps.
  \end{itemize}
\end{frame}

\begin{frame}[fragile]{Dessiner avec TikZ}
\begin{CodeRendu}[title={Un triangle rectangle}]
\begin{tikzpicture}[scale=0.55]
  \draw[thick] (0,0) -- (4,0)
        -- (4,3) -- cycle;
  \draw (3.6,0) |- (4,0.4);
  \node[below] at (2,0) {$4$};
  \node[right] at (4,1.5) {$3$};
  \node[above left] at (2,1.5) {$c$};
\end{tikzpicture}
\end{CodeRendu}
  \small
  \begin{itemize}
    \item On décrit la figure par des coordonnées : \texttt{(4,3)} est le
      point d'abscisse $4$ et d'ordonnée $3$.
    \item L'IA produit très bien du TikZ ; votre travail : ajuster une
      coordonnée, une étiquette, une échelle --- à la main.
  \end{itemize}
\end{frame}

% ═══════════════════════════════════════════════════════════════════════
\section{Travailler avec l'IA}
% ═══════════════════════════════════════════════════════════════════════

\begin{frame}{Le bon partage des rôles}
  \small
  \begin{columns}[T]
    \begin{column}{0.48\textwidth}
      \begin{Rem}{Ce que l'IA fait bien}
        \footnotesize
        \begin{itemize}
          \item Le squelette d'un document complet.
          \item La mise en forme : tableaux, listes, en-têtes.
          \item Les figures TikZ et les formules longues.
          \item Expliquer un message d'erreur.
          \item Proposer des variantes d'une question.
        \end{itemize}
      \end{Rem}
    \end{column}
    \begin{column}{0.48\textwidth}
      \begin{AuTableau}[Ce qui reste votre travail]
        \footnotesize
        \begin{itemize}
          \item Vérifier les mathématiques : calculs, corrigé, unités.
          \item Ajuster le niveau, le vocabulaire, le contexte québécois.
          \item Repérer l'anglais (Question, Points, Solution) et le
            point décimal.
          \item Refuser les extensions inutiles ou inventées.
          \item Comprendre chaque bloc avant de le garder.
        \end{itemize}
      \end{AuTableau}
    \end{column}
  \end{columns}
\end{frame}

\begin{frame}{L'anatomie d'un bon prompt}
  \small
  \begin{center}
  \begin{tabular}{lp{9.5cm}}
    \toprule
    Contexte & qui vous êtes, pour quel niveau, quelle situation \\
    Tâche & ce qu'il faut produire, et son contenu \\
    Contraintes techniques & classe, pdfLaTeX sur Overleaf, extensions
      permises, français, virgule décimale, longueur \\
    Format de sortie & un seul fichier complet et compilable, commenté \\
    Vérification & demander le corrigé, les étapes de calcul, la liste
      des extensions et leur rôle \\
    \bottomrule
  \end{tabular}
  \end{center}
  \vspace{4pt}
  Les crochets \champ{comme ceci} dans les prompts qui suivent sont à
  remplacer par vos propres choix.
\end{frame}

\begin{frame}{Prompt type : le squelette d'un document}
  \begin{Prompt}{Créer un document de départ}
    Tu es un enseignant de mathématiques \champ{au secondaire ou au
    collégial} au Québec et un expert de LaTeX.

    Écris un document LaTeX complet de classe article, qui compile sans
    erreur avec pdfLaTeX sur Overleaf. Il sera utilisé pour
    \champ{une fiche de notes sur les fonctions affines, 3\textsuperscript{e}
    secondaire} \champ{ou : un résumé sur les limites, Calcul différentiel au
    collégial}.

    Contraintes : texte en français ; extensions limitées à fontenc,
    babel (french), amsmath, amssymb, siunitx (virgule décimale), geometry
    et tikz ; aucune extension non standard ; pas de gras ni d'italique
    décoratifs.

    Ajoute un commentaire court à côté de chaque ligne du préambule pour
    expliquer son rôle. Donne uniquement le fichier .tex.
  \end{Prompt}
\end{frame}

\begin{frame}{Prompt type : corriger une erreur de compilation}
  \begin{Prompt}{Comprendre et réparer}
    Mon document LaTeX ne compile pas sur Overleaf (pdfLaTeX). Voici le
    message d'erreur complet :

    \champ{coller le message, par exemple : ! Missing \$ inserted.
    l.27 ...}

    Voici les lignes 20 à 30 de mon fichier :

    \champ{coller les lignes}

    Explique-moi en une ou deux phrases la cause de l'erreur, puis donne
    seulement les lignes corrigées. Ne réécris pas le reste du document.
  \end{Prompt}
  \small Donner le message exact et le passage concerné : sans eux,
  l'IA devine. Demander seulement les lignes corrigées : on garde le
  contrôle de son fichier.
\end{frame}

\begin{frame}{Prompts types : comprendre, modifier, dessiner}
  \begin{columns}[T]
    \begin{column}{0.49\textwidth}
      \begin{Prompt}{Expliquer un code}
        \scriptsize Voici un extrait de code LaTeX produit par une IA :
        \champ{extrait}. Explique ligne par ligne ce que fait chaque
        commande, en termes simples, pour une personne qui débute.
        Signale les lignes qui ne sont pas nécessaires.
      \end{Prompt}
      \vspace{2pt}
      \begin{Prompt}{Modification ciblée}
        \scriptsize Dans mon document, je veux \champ{ajouter une
        colonne « Points » au tableau de la section 2}. Donne seulement
        le bloc modifié, et indique où le placer.
      \end{Prompt}
    \end{column}
    \begin{column}{0.49\textwidth}
      \begin{Prompt}{Une figure TikZ}
        \scriptsize Écris le code TikZ d'une figure pour une fiche de
        \champ{3\textsuperscript{e} secondaire} : \champ{un triangle rectangle $ABC$,
        rectangle en $C$, avec $AC = 6$ cm et $BC = 8$ cm, le symbole
        d'angle droit et les mesures des côtés}. Uniquement
        l'environnement tikzpicture, avec des coordonnées entières et
        une échelle de \champ{0,5}, sans extension autre que tikz.
        Commente chaque ligne.
      \end{Prompt}
    \end{column}
  \end{columns}
\end{frame}

\begin{frame}{Avant de coller le code de l'IA}
  \begin{Rem}{Liste de vérification}
    \small
    \begin{enumerate}
      \item Je compile le code tel quel, avant toute modification.
      \item Je reconnais les trois parties : classe, préambule, corps.
      \item Chaque extension a un rôle que je peux nommer.
      \item Les nombres décimaux ont une virgule ; le texte est en
        français, y compris dans les en-têtes et les tableaux de points.
      \item J'ai refait moi-même les calculs du corrigé.
      \item Le niveau, le vocabulaire et le contexte conviennent à mes
        élèves ou à mes étudiants.
      \item Je sais où modifier un titre, une question, une formule, une
        coordonnée de figure.
    \end{enumerate}
  \end{Rem}
\end{frame}

% ═══════════════════════════════════════════════════════════════════════
\section{Livrable 1 : une activité d'enseignement}
% ═══════════════════════════════════════════════════════════════════════

\begin{frame}{La consigne}
  \begin{Exercice}{Présenter une activité mathématique, en 2 à 3 pages}
    \small
    \begin{itemize}
      \item Classe \texttt{article}, compilée sans erreur sur Overleaf.
      \item Contenu : titre ; niveau visé, durée, intention, préalables,
        matériel ; mise en situation ; déroulement ; tâche de l'élève ou de l'étudiant ;
        corrigé et pistes d'intervention.
      \item Au moins une figure (TikZ ou image) ou un tableau.
      \item Des mathématiques en mode mathématique : formules propres,
        virgule décimale, unités avec \cmd{qty}.
      \item Un en-tête de page avec le niveau et le titre de l'activité.
    \end{itemize}
  \end{Exercice}
  \small Point de départ possible : le fichier \texttt{02-activite.tex},
  ou un squelette produit par l'IA avec le prompt suivant.
\end{frame}

\begin{frame}{Prompt type : la fiche d'activité}
  \begin{Prompt}{Fiche de présentation d'une activité}
    \footnotesize Tu es conseiller pédagogique en mathématiques
    \champ{au secondaire ou au collégial} au Québec. Rédige en LaTeX (classe article, pdfLaTeX,
    Overleaf) la présentation d'une activité de \champ{75} minutes pour
    des \champ{élèves de 3\textsuperscript{e} secondaire, ou étudiants d'un
    cours collégial} sur \champ{la relation de
    Pythagore, dans le contexte de la taille des écrans de télévision}.

    Structure : titre ; un tableau de présentation (niveau, durée,
    intention, préalables, matériel) ; mise en situation ; déroulement
    minuté ; tâche de l'élève en \champ{4} questions graduées ; corrigé
    détaillé ; pistes d'intervention pour deux erreurs fréquentes.

    Contraintes : 2 à 3 pages ; français ; une figure TikZ simple ; un
    en-tête avec fancyhdr ; siunitx avec virgule décimale pour tous les
    nombres décimaux et les unités ; extensions standard seulement.
    Donne un seul fichier .tex complet, commenté.
  \end{Prompt}
\end{frame}

\begin{frame}{Un exemple de rendu : \texttt{02-activite.tex}}
  \begin{columns}[c]
    \begin{column}{0.33\textwidth}
      \centering
      \fbox{\includegraphics[height=0.66\textheight, page=1]{02-activite.pdf}}
    \end{column}
    \begin{column}{0.33\textwidth}
      \centering
      \fbox{\includegraphics[height=0.66\textheight, page=2]{02-activite.pdf}}
    \end{column}
    \begin{column}{0.3\textwidth}
      \small
      \begin{itemize}
        \item Titre, en-tête et pied de page.
        \item Tableau de présentation (liste \texttt{description}).
        \item Figure TikZ du téléviseur.
        \item Tâche et corrigé numérotés a), b), c), d).
      \end{itemize}
    \end{column}
  \end{columns}
\end{frame}

\begin{frame}[fragile]{Deux blocs de la fiche à savoir modifier}
  \begin{columns}[T]
    \begin{column}{0.5\textwidth}
\begin{CodeTeX}[title={L'en-tête de page (préambule)}]
\usepackage{fancyhdr}
\pagestyle{fancy}
\fancyhf{}                % on efface tout
\lhead{Mathématique, 3\textsuperscript{e} secondaire}
\rhead{Activité : Pythagore au salon}
\cfoot{\thepage}          % numéro au centre
\end{CodeTeX}
    \end{column}
    \begin{column}{0.47\textwidth}
\begin{CodeTeX}[title={Une question et sa réponse (corps)}]
\begin{enumerate}[label=\alph*)]
  \item Calcule la largeur, en
    centimètres, d'un écran de
    \qty{55}{po}.
\end{enumerate}
...
$k = \dfrac{55}{\sqrt{337}}
   \approx \num{2,996}$
\end{CodeTeX}
    \end{column}
  \end{columns}
  \vspace{4pt}
  \small \texttt{[label=\textbackslash alph*)]} (extension
  \texttt{enumitem}) numérote a), b), c). Changer le niveau, le titre ou
  une valeur numérique : une seule ligne à modifier, à la main.
\end{frame}

% ═══════════════════════════════════════════════════════════════════════
\section{Livrable 2 : un examen avec la classe exam}
% ═══════════════════════════════════════════════════════════════════════

\begin{frame}{Pourquoi la classe exam}
  \small
  \begin{itemize}
    \item Des questions, sous-questions et choix multiples numérotés
      automatiquement.
    \item Des points par question ; le total et la grille de correction se
      calculent seuls.
    \item Les solutions écrites sous chaque question : un seul fichier pour
      la copie de l'élève et le corrigé.
    \item Des espaces de réponse réglables : vide, lignes, cadre.
    \item Des en-têtes et pieds de page avec « Page 1 de 3 ».
  \end{itemize}
  \vspace{4pt}
  \begin{Rem}{Un seul interrupteur}
    \small Avec \cmd{printanswers} dans le préambule, les solutions
    s'affichent ; en le mettant en commentaire (\texttt{\%}), elles
    disparaissent et l'espace de réponse est conservé.
  \end{Rem}
\end{frame}

\begin{frame}[fragile]{La structure d'un examen}
  \begin{columns}[T]
    \begin{column}{0.56\textwidth}
\begin{CodeTeX}[title={Le corps de 03-examen.tex (extrait)}]
\begin{questions}
\question[2] Lequel de ces triplets ... ?
\begin{choices}
  \CorrectChoice $3$, $4$ et $5$
  \choice $4$, $5$ et $6$
\end{choices}
\question Un triangle rectangle ...
\begin{parts}
  \part[2] Calcule l'hypoténuse.
  \begin{solution}
    $c = \sqrt{6^2 + 8^2} = \qty{10}{\centi\metre}$.
  \end{solution}
\end{parts}
\end{questions}
\end{CodeTeX}
    \end{column}
    \begin{column}{0.41\textwidth}
      \footnotesize
      \begin{itemize}
        \item \texttt{questions} : la liste de toutes les questions.
        \item \cmd{question[2]} : une question de 2 points.
        \item \texttt{parts}, \cmd{part[2]} : les sous-questions a), b).
        \item \texttt{choices}, \cmd{choice} : un choix multiple ;
          \cmd{CorrectChoice} marque la bonne réponse dans le corrigé.
        \item \texttt{solution} : la réponse, visible seulement avec
          \cmd{printanswers}.
      \end{itemize}
    \end{column}
  \end{columns}
\end{frame}

\begin{frame}[fragile]{Mettre la classe exam en français}
  \begin{columns}[T]
    \begin{column}{0.56\textwidth}
\begin{CodeTeX}[title={Préambule de 03-examen.tex (extrait)}]
\documentclass[11pt, addpoints]{exam}
% \printanswers   % corrigé : enlever le %
\pointpoints{point}{points}
\hqword{Question}
\hpword{Points}
\hsword{Résultat}
\htword{Total}
\renewcommand{\solutiontitle}{\noindent Solution : }
\pagestyle{headandfoot}
\firstpageheader{Mathématique, 3\textsuperscript{e} secondaire}
                {}{Examen --- Relation de Pythagore}
\firstpagefooter{}{Page \thepage\ de \numpages}{}
\end{CodeTeX}
    \end{column}
    \begin{column}{0.41\textwidth}
      \footnotesize
      \begin{itemize}
        \item \texttt{addpoints} : active le calcul du total.
        \item Sans ces réglages, la grille affiche Question, Points,
          Score : l'IA les oublie souvent.
        \item En-tête et pied de page : trois cases (gauche, centre,
          droite), entre accolades.
      \end{itemize}
    \end{column}
  \end{columns}
\end{frame}

\begin{frame}[fragile]{Grille, total et espaces de réponse}
\begin{CodeTeX}[title={Quelques commandes utiles}]
Cet examen compte \numquestions\ questions, pour un total de \numpoints\ points.
\gradetable[h][questions]          % grille de correction horizontale

\vspace{\stretch{1}}               % un espace qui partage la page restante
\begin{solution}[3cm]              % 3 cm réservés si le corrigé est caché
\fillwithlines{3cm}                % 3 cm de lignes pour écrire
\makeemptybox{3cm}                 % un cadre vide de 3 cm
\end{CodeTeX}
  \small Comme pour les sections, rien n'est numéroté ni additionné à la
  main : ajouter une question met à jour la numérotation, le total et la
  grille à la compilation suivante.
\end{frame}

\begin{frame}{Le rendu : copie de l'élève et corrigé}
  \begin{columns}[c]
    \begin{column}{0.36\textwidth}
      \centering
      \fbox{\includegraphics[height=0.7\textheight]{03-examen.pdf}}\\[2pt]
      \scriptsize sans \cmd{printanswers}
    \end{column}
    \begin{column}{0.36\textwidth}
      \centering
      \fbox{\includegraphics[height=0.7\textheight]{03-examen-corrige.pdf}}\\[2pt]
      \scriptsize avec \cmd{printanswers}
    \end{column}
    \begin{column}{0.25\textwidth}
      \small Le même fichier, une seule ligne de différence. La bonne
      réponse du choix multiple apparaît en gras dans le corrigé (réglage
      par défaut de la classe).
    \end{column}
  \end{columns}
\end{frame}

\begin{frame}{La consigne}
  \begin{Exercice}{Un examen sur le même thème que votre activité}
    \small
    \begin{itemize}
      \item Classe \texttt{exam}, avec l'option \texttt{addpoints} et la
        classe réglée en français.
      \item Au moins quatre questions, dont un choix multiple et une
        question à sous-questions ; des points pour chacune.
      \item Une grille de correction et le total calculés automatiquement.
      \item Une solution pour chaque question ; remettre deux PDF : la
        copie de l'élève et le corrigé.
      \item En-tête avec le niveau, pied de page « Page x de y ».
    \end{itemize}
  \end{Exercice}
\end{frame}

\begin{frame}{Prompt type : l'examen}
  \begin{Prompt}{Un examen avec la classe exam}
    \footnotesize Rédige en LaTeX un examen de \champ{3\textsuperscript{e}
    secondaire, ou d'un cours collégial} sur \champ{la relation de Pythagore}, avec la classe exam
    (option addpoints), compilable avec pdfLaTeX sur Overleaf.

    Contenu : \champ{5} questions pour un total de \champ{20} points ;
    au moins un choix multiple (choices, CorrectChoice), une question avec
    sous-questions (parts), et un problème en contexte réel. Mets une
    solution détaillée dans un environnement solution sous chaque question,
    et laisse un espace de réponse adapté.

    Mise en forme : la classe exam entièrement en français (pointpoints,
    hqword, hpword, hsword, htword, solutiontitle) ; une grille de
    correction horizontale ; en-tête avec le niveau et le titre ; pied de
    page « Page x de y ». Laisse \textbackslash printanswers en
    commentaire au début du fichier. Utilise siunitx avec virgule décimale
    pour les nombres et les unités. Donne un seul fichier .tex.
  \end{Prompt}
\end{frame}

\begin{frame}{Prompts types : réutiliser ses questions}
  \begin{columns}[T]
    \begin{column}{0.49\textwidth}
      \begin{Prompt}{Une version B}
        \scriptsize Voici mon examen en LaTeX : \champ{fichier}. Produis
        une version B de même difficulté : mêmes questions, mêmes points,
        mais d'autres valeurs numériques, choisies pour que les réponses
        restent simples. Recalcule toutes les solutions et montre-moi les
        calculs pour que je les vérifie.
      \end{Prompt}
    \end{column}
    \begin{column}{0.49\textwidth}
      \begin{Prompt}{Ajouter une question}
        \scriptsize Ajoute à mon examen une question de \champ{4} points
        sur \champ{la réciproque de la relation de Pythagore}, avec deux
        sous-questions et leur solution. Donne seulement le bloc
        \textbackslash question à insérer, dans le même style que les
        autres.
      \end{Prompt}
    \end{column}
  \end{columns}
  \vspace{6pt}
  \small Dans les deux cas : refaire soi-même les calculs du corrigé avant
  de distribuer la copie.
\end{frame}

% ═══════════════════════════════════════════════════════════════════════
\section{Remise et ressources}
% ═══════════════════════════════════════════════════════════════════════

\begin{frame}{Avant de remettre}
  \begin{Rem}{Liste de vérification finale}
    \small
    \begin{enumerate}
      \item Les deux documents compilent sans erreur (aucun message rouge).
      \item L'activité fait 2 à 3 pages et contient toutes les rubriques
        demandées, avec au moins une figure ou un tableau.
      \item L'examen affiche la grille et le bon total ; le corrigé
        s'obtient en activant \cmd{printanswers}.
      \item Tout est en français ; les décimaux ont une virgule.
      \item Les mathématiques ont été vérifiées par vous, pas seulement
        par l'IA.
      \item Vous pouvez expliquer à quoi sert chaque ligne du préambule.
    \end{enumerate}
    À remettre : \texttt{activite.pdf}, \texttt{examen.pdf},
    \texttt{examen-corrige.pdf}, et les fichiers \texttt{.tex}.
  \end{Rem}
\end{frame}

\begin{frame}{Pour aller plus loin}
  \small
  \begin{itemize}
    \item Overleaf, Learn LaTeX in 30 minutes et la documentation en ligne :
      \texttt{overleaf.com/learn}.
    \item Detexify, pour trouver un symbole en le dessinant :
      \texttt{detexify.kirelabs.org}.
    \item La documentation de la classe exam, très complète, avec de
      nombreux exemples : \texttt{ctan.org/pkg/exam}.
    \item La documentation de siunitx : \texttt{ctan.org/pkg/siunitx}.
    \item Les trois fichiers d'exemple de ce laboratoire :
      \texttt{01-minimal.tex}, \texttt{02-activite.tex},
      \texttt{03-examen.tex}.
  \end{itemize}
  \vspace{6pt}
  \begin{Classe}{Pour votre pratique}
    \small Commencez une banque de questions dès vos premiers stages : un fichier
    \texttt{exam} par notion, que vous enrichirez et réutiliserez pendant
    toute votre carrière.
  \end{Classe}
\end{frame}

\appendix

\begin{frame}[fragile]{Solutions de l'exercice « Reproduire ces formules »}
\begin{CodeTeX}[title={Une solution possible}]
% 1.
$f(x) = -2(x - 3)^2 + 1$
% 2.
L'aire d'un disque de rayon \qty{5}{\centi\metre} est
$A = \pi r^2 \approx \qty{78,54}{\centi\metre\squared}$.
% 3.
\[ x = \frac{-b \pm \sqrt{b^2 - 4ac}}{2a} \]
% 4.
\begin{align*}
  2(x + 4) &= 18 \\
  x + 4 &= 9 \\
  x &= 5
\end{align*}
% 5.
$\sin(\theta) = \dfrac{\text{opposé}}{\text{hypoténuse}}$
\end{CodeTeX}
\end{frame}

\begin{frame}[fragile]{Solution de l'exercice « Réparer »}
  \begin{columns}[T]
    \begin{column}{0.48\textwidth}
\begin{CodeTeX}[title={Code corrigé}]
\section{Exercice}
Calculer $x^2$ si $x = 3$.
\begin{itemize}
  \item La fraction $\frac{1}{2}$.
  \item Le rabais : 20\%.
\end{itemize}
\end{CodeTeX}
    \end{column}
    \begin{column}{0.48\textwidth}
      \small
      \begin{enumerate}
        \item Accolade fermante oubliée après le titre de section.
        \item \texttt{x\^{}2} hors du mode mathématique.
        \item Accolade fermante oubliée dans \cmd{frac}.
        \item \texttt{\%} non protégé : la fin de la ligne devient un
          commentaire.
        \item \cmd{begin\{itemize\}} fermé par
          \cmd{end\{enumerate\}}.
      \end{enumerate}
    \end{column}
  \end{columns}
\end{frame}


\end{document}
